\documentclass[11pt,a4paper]{article}

\usepackage[T1]{fontenc}
\usepackage[utf8]{inputenc}
\usepackage{lmodern}
\usepackage{microtype}
\usepackage{geometry}
\geometry{margin=30mm}
\usepackage{amsmath,amssymb,amsthm,mathtools}
\usepackage{enumitem}
\usepackage{booktabs}
\usepackage{xcolor}
\usepackage[colorlinks=true,linkcolor=blue!45!black,
  citecolor=blue!45!black,urlcolor=blue!45!black]{hyperref}
\hypersetup{
  pdftitle={Dihedral-Angle Sums in Orthocentric, Power-k, and Isodynamic Tetrahedra},
  pdfauthor={Ismail Hammoudeh},
  pdfsubject={Sharp ranges for orthocentric, Power-k, and isodynamic tetrahedra},
  pdfkeywords={Power-k tetrahedron, orthocentric tetrahedron, 4-ball tetrahedron, balloon tetrahedron, isodynamic tetrahedron, dihedral angle, solid angle, kite}
}

\newtheorem{theorem}{Theorem}[section]
\newtheorem{lemma}[theorem]{Lemma}
\newtheorem{proposition}[theorem]{Proposition}
\newtheorem{corollary}[theorem]{Corollary}
\newtheorem{conjecture}[theorem]{Conjecture}
\theoremstyle{definition}
\newtheorem{definition}[theorem]{Definition}
\theoremstyle{remark}
\newtheorem{remark}[theorem]{Remark}

\newcommand{\R}{\mathbb{R}}
\newcommand{\vol}{\operatorname{vol}}
\newcommand{\atanTwo}{\operatorname{atan2}}

\title{Dihedral-Angle Sums in Orthocentric, Power-$k$,\\
and Isodynamic Tetrahedra}
\author{Ismail Hammoudeh}
\date{}

\begin{document}
\maketitle

\begin{abstract}
Let $\Sigma(T)$ denote the sum of the six internal dihedral angles of a
nondegenerate Euclidean tetrahedron $T$.  We organize several special
families around the \emph{Power-$k$ family}
\[
 d_{ij}^{\,k}=y_i+y_j,\qquad k\ne0,
\]
with positive pair sums and Euclidean realizability.  The case $k=2$ is
precisely the full signed orthocentric family.  We prove that every
orthocentric tetrahedron satisfies
\[
  \Sigma(T)\geq 6\arccos\frac13,
\]
with equality exactly for the regular tetrahedron, and that its exact range
is
\[
 \left[6\arccos\frac13,\,3\pi\right).
\]
The positive branch of the Power-$2$ family has the smaller exact range
\[
 [6\arccos(1/3),5\pi/2).
\]
We prove that for every $k<1$, $k\ne0$, and every $k>2$, the full Power-$k$
family instead has the universal exact range $(2\pi,3\pi)$.  For
$1<k<2$ the global lower bound remains open.  We prove that the regular
tetrahedron is a strict local minimum and reduce the conjectured global
bound to an exponent-monotonicity inequality; Schoenberg's theorem ensures
that every required intermediate metric is a nondegenerate Euclidean
tetrahedron.
The orthocentric proof combines the constant polar-measure sum with a new
convex lower envelope for ordinary solid-angle content.  In the normalized
positive power-mean parametrization, the limit $k\to0$ is the isodynamic
metric; we prove that its exact range is again $(2\pi,3\pi)$.  Finally, as a
matter of logical status, the regular lower bound recently stated for the
$k=1$ (4-ball or balloon) family is not used as an input.
\end{abstract}

\noindent\textbf{Keywords.}
Power-$k$ tetrahedron; orthocentric tetrahedron; 4-ball tetrahedron; balloon
tetrahedron; isodynamic tetrahedron; isogonic tetrahedron; dihedral angle;
solid angle; polar sine; kite.

\noindent\textbf{2020 Mathematics Subject Classification.}
Primary 51M04; Secondary 52B10, 52B11.

\section{Introduction}

For a nondegenerate Euclidean tetrahedron $T$, write $\Sigma(T)$ for the sum
of its six internal dihedral angles.  The sharp universal bounds
\begin{equation}\label{eq:gaddum}
  2\pi<\Sigma(T)<3\pi
\end{equation}
were proved by Jarden (Juzuk) in 1941; see \cite{Jarden1941} and the later
short treatment of Gaddum \cite{Gaddum1952}.  Neither endpoint is attained by
a nondegenerate tetrahedron.

The regular tetrahedron has six equal dihedral angles $\arccos(1/3)$ and
therefore
\begin{equation}\label{eq:regular-value}
  \Sigma_{\rm reg}=6\arccos\frac13
  =2.350959312\ldots\,\pi.
\end{equation}
The purpose of this article is to determine when this regular value remains
the lower endpoint after one restricts the metric to natural tetrahedral
families, and when the infimum falls all the way to the universal value
$2\pi$.

We use the terminology \emph{Power-$k$ tetrahedron} for the family described
by
\begin{equation}\label{eq:power-k-intro}
 d_{ij}^{\,k}=y_i+y_j,\qquad k\in\mathbb R\setminus\{0\},
\end{equation}
subject to positive pair sums and Euclidean realizability.  This notation
exposes two important members of the same family.  At $k=2$ one obtains the
signed additive metric
$d_{ij}^2=y_i+y_j$, which is exactly the orthocentric family
\cite{EdmondsHajjaMartini2005,HajjaHammoudeh2014}.  At $k=1$ the metric is
edge-additive.  The latter family is the classical 4-ball or balloon family;
its equivalent characterizations are discussed, for example, by Katsuura
\cite{Katsuura2021}.  Korotov and K\v{r}\'\i\v{z}ek recently stated that its
dihedral-angle sum has the exact range
$[6\arccos(1/3),3\pi)$ \cite[Theorem~9]{KorotovKrizek2026}.  The proof in
that preprint gives the symmetric part of the required global maximization
and refers to the nonsymmetric case as a calculation without displaying it.
We therefore do not use the $k=1$ lower bound as an established input.

Our main results and reductions are as follows.
\begin{enumerate}[label=\textup{(\roman*)},leftmargin=8mm]
\item We prove the regular lower bound for the full orthocentric, hence
Power-$2$, family, including the component with one negative additive
parameter.  Equilateral-base kites then show that the exact range is
$[6\arccos(1/3),3\pi)$.
\item For the positive Power-$2$ branch we combine the new lower bound with
the sharp nonobtuse upper bound to obtain the exact range
$[6\arccos(1/3),5\pi/2)$.
\item For every $k<1$, $k\ne0$, and every $k>2$, we construct convex-planar
degenerations inside the full Power-$k$ family and prove that the exact range
is the full universal interval $(2\pi,3\pi)$.  The interval $1<k<2$
remains open.  We prove strict local minimality of the regular tetrahedron
throughout this interval and give a rigorous reduction to exponent
monotonicity using Schoenberg's snowflake theorem.
\item In the normalized positive power-mean parametrization, the limit
$k\to0$ is the isodynamic metric.  We determine its exact range to be
$(2\pi,3\pi)$ by two explicit endpoint paths and continuity.
\end{enumerate}

The orthocentric lower bound is the conceptual core of the paper.  Hajja and
Hammoudeh proved that the four polar measures of the corner angles of every
orthocentric tetrahedron add to $\pi$ \cite[Theorem~1(b)]{HajjaHammoudeh2014}.
We combine that identity with an elementary comparison between polar measure
and ordinary solid-angle content.  The minimizing acute corner at fixed polar
measure is equiangular, and its content is a strictly convex function of the
polar measure.  Jensen's inequality then gives the regular lower bound in the
all-acute case, while a separate estimate handles the possible right or
obtuse corner.

The following table summarizes the status proved or used in this paper.  The
entry at $k=0$ is a normalized limiting family, not a literal substitution in
\eqref{eq:power-k-intro}.
\begin{center}
\renewcommand{\arraystretch}{1.18}
\begin{tabular}{ccc}
\toprule
parameter & family & exact range of $\Sigma$ \\
\midrule
$k<1$, $k\ne0$ & full $\mathcal P_k$ & $(2\pi,3\pi)$ \\
$k=1$ & 4-ball/balloon $\mathcal P_1$ & $[\Sigma_{\rm reg},3\pi)$ stated in \cite{KorotovKrizek2026} \\
$1<k<2$ & full $\mathcal P_k$ & $[\Sigma_{\rm reg},3\pi)$ conjectured \\
$k=2$ & orthocentric $\mathcal P_2$ & $[\Sigma_{\rm reg},3\pi)$ \\
$k=2$ & positive branch $\mathcal P_2^+$ & $[\Sigma_{\rm reg},5\pi/2)$ \\
$k>2$ & full $\mathcal P_k$ & $(2\pi,3\pi)$ \\
$k\to0$ & normalized positive limit: isodynamic & $(2\pi,3\pi)$
\\
\bottomrule
\end{tabular}
\end{center}

The organization is as follows.  Section~\ref{sec:preliminaries} introduces
the two angle measures, the Power-$k$ terminology and its special cases, and
the constant polar-sum theorem.  Section~\ref{sec:acute-corner} proves the
local comparison and its crucial convexity.  Section~\ref{sec:main-proof}
proves the orthocentric lower bound, treating the all-acute and non-acute
cases separately.  Section~\ref{sec:kites} uses kites to obtain the exact
orthocentric range and the range of the positive Power-$2$ branch.
Section~\ref{sec:power-family} proves the exterior Power-$k$ ranges and gives
rigorous partial results for $1<k<2$.  Section~\ref{sec:isodynamic} proves the
exact isodynamic range.  The final section puts the different proof
mechanisms in perspective.

\section{Preliminaries}\label{sec:preliminaries}

\subsection{Content and the dihedral sum}

Consider a trihedral angle with vertex at the origin and unit vectors
$u_1,u_2,u_3$ along its three arms.  Set
\[
 a=u_2\mathbin{\cdot}u_3,\qquad
 b=u_3\mathbin{\cdot}u_1,\qquad
 c=u_1\mathbin{\cdot}u_2.
\]
Thus $a,b,c$ are the cosines of its three planar subangles.  Its polar sine is
\begin{equation}\label{eq:polar-sine}
 P=\sqrt{1-a^2-b^2-c^2+2abc}>0.
\end{equation}
The ordinary solid-angle content $\Omega\in(0,2\pi)$ is the area of the
spherical triangle cut out by the three arms on the unit sphere.  The standard
determinant formula is
\begin{equation}\label{eq:content-formula}
 \Omega=2\atanTwo\bigl(P,1+a+b+c\bigr),
\end{equation}
where $\atanTwo(y,x)\in(0,\pi)$ denotes the argument of $x+iy$ when $y>0$.
This notation keeps the correct branch when $1+a+b+c$ is negative.

Let $\Omega_1,\ldots,\Omega_4$ be the contents at the four vertices of a
tetrahedron, and let $\theta_1,\ldots,\theta_6$ be its internal dihedral
angles.  At each vertex, spherical excess gives
\[
 \Omega_i=(\text{sum of the three incident dihedral angles})-\pi.
\]
Summing over the four vertices counts each dihedral angle twice.  Hence
\begin{equation}\label{eq:content-dihedral}
 \sum_{i=1}^4\Omega_i=2\Sigma(T)-4\pi,
 \qquad
 \Sigma(T)=2\pi+\frac12\sum_{i=1}^4\Omega_i.
\end{equation}

\subsection{Pure corners and polar measure}

A corner is called \emph{acute}, \emph{right}, or \emph{obtuse} when its
three planar subangles are respectively all acute, all right, or all obtuse.
These are the pure corners relevant here.  Following
\cite{HajjaHammoudeh2014}, define their polar measure by
\begin{equation}\label{eq:polar-measure}
 \beta=
 \begin{cases}
  \arcsin P, & \text{for an acute corner},\\[1mm]
  \pi/2, & \text{for a right corner},\\[1mm]
  \pi-\arcsin P, & \text{for an obtuse corner}.
 \end{cases}
\end{equation}
In every case $\sin\beta=P$.

An orthocentric tetrahedron admits real additive parameters
$a_1,\ldots,a_4$ such that
\begin{equation}\label{eq:orthocentric-metric}
 d_{ij}^2=a_i+a_j\qquad(i\ne j),
\end{equation}
and all pair sums are positive; see
\cite{EdmondsHajjaMartini2005,HajjaHammoudeh2014}.  Consequently at most one
$a_i$ can be nonpositive.  The cosine rule shows that all three planar
subangles at vertex $i$ have the sign of $a_i$.  Thus every corner is pure,
at least three corners are acute, and at most one is right or obtuse.

\begin{definition}[Power-$k$ tetrahedron]\label{def:power-k}
Let $k\in\R\setminus\{0\}$.  A nondegenerate Euclidean tetrahedron is a
\emph{Power-$k$ tetrahedron} if its edge lengths satisfy
\begin{equation}\label{eq:power-k-metric}
 d_{ij}^{\,k}=y_i+y_j>0\qquad(i\ne j)
\end{equation}
for real parameters $y_1,\ldots,y_4$.  We denote the family by
$\mathcal P_k$.  Its \emph{positive branch} $\mathcal P_k^+$ consists of the
members admitting $y_i>0$ for all $i$.
\end{definition}

The requirement that all pair sums in \eqref{eq:power-k-metric} be positive
already implies that at most one $y_i$ is negative.  Multiplying every $y_i$
by the same positive constant changes only the scale of the tetrahedron, so
the Power-$k$ family has three shape parameters.

\begin{proposition}[Intrinsic Power-$k$ relation]\label{prop:power-opposite}
For fixed $k\ne0$, six positive edge lengths belong to the Power-$k$ family
if and only if
\begin{equation}\label{eq:power-opposite}
 d_{12}^k+d_{34}^k
 =d_{13}^k+d_{24}^k
 =d_{14}^k+d_{23}^k.
\end{equation}
Thus the definition is intrinsic and does not depend on a chosen set of
parameters.
\end{proposition}

\begin{proof}
If \eqref{eq:power-k-metric} holds, each expression in
\eqref{eq:power-opposite} equals $y_1+y_2+y_3+y_4$.  Conversely, put
$A_{ij}=d_{ij}^k$ and assume the three opposite-pair sums are equal.  Define
\[
 y_1=\frac{A_{12}+A_{13}-A_{23}}2,\quad
 y_2=\frac{A_{12}+A_{23}-A_{13}}2,\quad
 y_3=\frac{A_{13}+A_{23}-A_{12}}2,
\]
and $y_4=A_{14}-y_1$.  Then $y_i+y_j=A_{ij}$ for the three edges among
$1,2,3$.  The two remaining equalities follow from
$A_{14}+A_{23}=A_{13}+A_{24}=A_{12}+A_{34}$, proving
\eqref{eq:power-k-metric} for every edge.
\end{proof}

The special exponents $1$ and $2$ deserve to be separated carefully.
For $k=2$, \eqref{eq:power-k-metric} is exactly
\eqref{eq:orthocentric-metric}; hence
\begin{equation}\label{eq:p2-orthocentric}
 \mathcal P_2=\{\text{orthocentric tetrahedra}\}.
\end{equation}
The positive branch $\mathcal P_2^+$ is the interior of the nonobtuse
component of the orthocentric family; allowing zero parameters gives its
right-corner boundary.  Indeed, writing $y_i=x_i^2$ gives
\begin{equation}\label{eq:k2-family}
 d_{ij}^2=x_i^2+x_j^2,
 \qquad x_i>0,
\end{equation}
whereas an orthocentric tetrahedron with one pure obtuse corner requires one
negative additive parameter.

For $k=1$, the face inequalities remove the signed component automatically.
If $d_{ij}=y_i+y_j$, then in any face containing vertex $i$,
\[
 d_{ij}+d_{ik}>d_{jk}
 \quad\Longleftrightarrow\quad 2y_i>0.
\]
Hence every realizable Power-$1$ tetrahedron has $y_i>0$ for all $i$.  Thus
$\mathcal P_1=\mathcal P_1^+$ and this family is exactly the classical
edge-additive, balloon, or 4-ball family; see \cite{Katsuura2021} for these
equivalent viewpoints.  The logical status of its proposed sharp range is
discussed in Section~\ref{sec:power-family}.

Finally, the positive branch has a natural normalized formulation.  For
$x_i>0$ set
\begin{equation}\label{eq:power-mean-family}
 D_{ij}^{(k)}=\left(\frac{x_i^k+x_j^k}{2}\right)^{1/k}.
\end{equation}
For each fixed $k\ne0$ this parametrizes the same shapes as
$\mathcal P_k^+$.  If instead the $x_i$ are held fixed while $k\to0$, then
\begin{equation}\label{eq:k0-limit}
 \lim_{k\to0}D_{ij}^{(k)}=\sqrt{x_ix_j}.
\end{equation}
Writing $p_i=\sqrt{x_i}$ gives $d_{ij}=p_ip_j$, the isodynamic metric.  Thus
the isodynamic family is the normalized positive $k\to0$ limit, not a
literal $k=0$ instance of \eqref{eq:power-k-metric}.

We use the following established result as an external input.

\begin{theorem}[Constant polar-measure sum]\label{thm:polar-sum}
For every nondegenerate orthocentric tetrahedron, the polar measures of its
four corner angles satisfy
\begin{equation}\label{eq:polar-sum}
 \beta_1+\beta_2+\beta_3+\beta_4=\pi.
\end{equation}
\end{theorem}

\begin{proof}
This is \cite[Theorem~1(b), Section~5]{HajjaHammoudeh2014}.  The branch
choice in \eqref{eq:polar-measure} is essential in the obtuse case.
\end{proof}

\section{The least content of an acute corner}\label{sec:acute-corner}

The key local problem is the following: if the polar measure of an acute
corner is fixed, how small can its ordinary content be?  The minimizing
corner is equiangular.

Define
\begin{equation}\label{eq:phi}
 \phi(q)=(1-q)^2(1+2q)=1-3q^2+2q^3,
 \qquad 0\leq q\leq1.
\end{equation}
The function $\phi$ decreases strictly from $1$ to $0$.  Given
$\beta\in[0,\pi/2]$, let $q=q(\beta)\in[0,1]$ be the unique solution of
\begin{equation}\label{eq:q-beta}
 \sin^2\beta=\phi(q).
\end{equation}
Define
\begin{equation}\label{eq:def-f}
 f(\beta)=3\arccos\frac{q(\beta)}{1+q(\beta)}-\pi.
\end{equation}
Geometrically, $f(\beta)$ is the content of the equiangular trihedral angle
whose three planar subangles have cosine $q(\beta)$.

\begin{lemma}[Acute-corner comparison]\label{lem:acute-corner}
If an acute trihedral angle has polar measure $\beta$ and content $\Omega$,
then
\begin{equation}\label{eq:acute-comparison}
 \Omega\geq f(\beta).
\end{equation}
Equality holds if and only if its three planar subangles are equal.
\end{lemma}

\begin{proof}
Use the notation $a,b,c>0$ from \eqref{eq:polar-sine}, and put
$s=a+b+c$.  The elementary inequalities
\[
 a^2+b^2+c^2\geq\frac{s^2}{3},
 \qquad
 abc\leq\left(\frac{s}{3}\right)^3
\]
give
\begin{align*}
 P^2
 &=1-a^2-b^2-c^2+2abc\\
 &\leq1-\frac{s^2}{3}+\frac{2s^3}{27}
 =\phi\left(\frac{s}{3}\right).
\end{align*}
By \eqref{eq:q-beta}, $P^2=\phi(q(\beta))$.  Since $\phi$ is decreasing,
\[
 q(\beta)\geq\frac{s}{3},
 \qquad
 1+s\leq1+3q(\beta).
\]
For fixed $P>0$, the expression $2\atanTwo(P,X)$ decreases as $X$
increases.  Formula \eqref{eq:content-formula} therefore yields
\begin{equation}\label{eq:f-atan}
 \Omega
 \geq
 2\arctan\frac{\sin\beta}{1+3q(\beta)}.
\end{equation}

For an equiangular corner with $a=b=c=q$, the spherical cosine rule shows
that each angle $\delta$ of the corresponding equilateral spherical triangle
satisfies
\[
 \cos\delta=\frac{q}{1+q}.
\]
Its spherical excess is $3\delta-\pi$, so the right-hand side of
\eqref{eq:f-atan} is precisely $f(\beta)$.  Equality in the two elementary
inequalities above occurs simultaneously exactly when $a=b=c$.
\end{proof}

\begin{lemma}[Convexity]\label{lem:convexity}
The function $f$ is strictly increasing and strictly convex on
$[0,\pi/2]$.  Moreover,
\begin{equation}\label{eq:linear-bound}
 f(0)=0,\qquad f'_+(0)=\frac12,\qquad
 f(\beta)\geq\frac{\beta}{2}.
\end{equation}
\end{lemma}

\begin{proof}
It is convenient to regard both $\beta$ and $f$ as functions of
$q\in[0,1]$:
\[
 \beta(q)=\arcsin\bigl((1-q)\sqrt{1+2q}\bigr),
 \qquad
 f(q)=3\arccos\frac{q}{1+q}-\pi.
\]
For $0<q<1$, direct differentiation gives
\begin{equation}\label{eq:parameter-derivatives}
 \frac{d\beta}{dq}
 =-\frac{3}{\sqrt{(1+2q)(3-2q)}},
 \qquad
 \frac{df}{dq}
 =-\frac{3}{(1+q)\sqrt{1+2q}}.
\end{equation}
Consequently,
\begin{equation}\label{eq:f-prime}
 f'(\beta)=\frac{\sqrt{3-2q}}{1+q}>0.
\end{equation}
Differentiating once more and using \eqref{eq:parameter-derivatives},
\begin{equation}\label{eq:f-second}
 f''(\beta)
 =\frac{(4-q)\sqrt{1+2q}}{3(1+q)^2}>0.
\end{equation}
Thus $f$ is strictly increasing and strictly convex.  As $\beta\downarrow0$
we have $q\uparrow1$, and \eqref{eq:f-prime} gives $f'_+(0)=1/2$.
The graph of a convex function lies above every supporting tangent, which
proves $f(\beta)\geq\beta/2$.
\end{proof}

At the regular value $\beta=\pi/4$, equation \eqref{eq:q-beta} gives
$q=1/2$, and hence
\begin{equation}\label{eq:f-regular}
 f\left(\frac{\pi}{4}\right)
 =3\arccos\frac13-\pi.
\end{equation}

\section{The lower bound}\label{sec:main-proof}

We first treat orthocentric tetrahedra whose four corners are acute.

\begin{proposition}[All-acute case]\label{prop:acute-case}
If all four corners of an orthocentric tetrahedron $T$ are acute, then
\begin{equation}\label{eq:acute-main}
 \Sigma(T)\geq6\arccos\frac13.
\end{equation}
Equality holds if and only if $T$ is regular.
\end{proposition}

\begin{proof}
By Lemma~\ref{lem:acute-corner}, Theorem~\ref{thm:polar-sum}, Jensen's
inequality, and \eqref{eq:f-regular},
\begin{align*}
 \sum_{i=1}^4\Omega_i
 &\geq\sum_{i=1}^4 f(\beta_i)\\
 &\geq4f\left(\frac{\beta_1+\beta_2+\beta_3+\beta_4}{4}\right)\\
 &=4f\left(\frac{\pi}{4}\right)\\
 &=12\arccos\frac13-4\pi.
\end{align*}
Equation \eqref{eq:content-dihedral} now gives \eqref{eq:acute-main}.

Suppose equality holds.  Strict convexity forces
$\beta_1=\cdots=\beta_4=\pi/4$, and equality in
Lemma~\ref{lem:acute-corner} forces each corner to be equiangular.  In
particular, at any vertex the three planar subangles have cosine $1/2$.
Using \eqref{eq:orthocentric-metric} at, say, vertex $4$, their cosines are
\[
 \frac{a_4}{\sqrt{(a_4+a_i)(a_4+a_j)}}
 \qquad (1\leq i<j\leq3).
\]
Their equality implies $a_1=a_2=a_3$, and their common value $1/2$ then
implies $a_4=a_1$.  Thus all $a_i$ are equal and all six edges are equal.
Conversely, the regular tetrahedron plainly gives equality.
\end{proof}

It remains to ensure that the possible right or obtuse corner cannot lower
the sum.  In fact it leads to a stronger estimate.

\begin{lemma}[Obtuse-corner comparison]\label{lem:obtuse-corner}
If a pure obtuse corner has polar measure $\beta\in(\pi/2,\pi)$ and content
$\Omega$, then
\begin{equation}\label{eq:obtuse-comparison}
 \Omega\geq\beta.
\end{equation}
The same inequality holds with equality for a right corner, where
$\Omega=\beta=\pi/2$.
\end{lemma}

\begin{proof}
For an obtuse corner write
\[
 a=-A,\qquad b=-B,\qquad c=-C,
 \qquad 0<A,B,C<1,
\]
and set $r=A+B+C$.  By \eqref{eq:polar-sine},
\[
 1-P^2=A^2+B^2+C^2+2ABC.
\]
Since $AB+AC+BC\geq ABC$,
\[
 r^2=A^2+B^2+C^2+2(AB+AC+BC)\geq1-P^2.
\]
Because $\beta=\pi-\arcsin P$, this says
\[
 r\geq\sqrt{1-P^2}=-\cos\beta.
\]
Thus $1-r\leq1+\cos\beta$.  Using the monotonicity of
$\atanTwo(P,X)$ in its second argument,
\begin{align*}
 \Omega
 &=2\atanTwo(P,1-r)\\
 &\geq2\atanTwo(\sin\beta,1+\cos\beta)
 =\beta.
\end{align*}
For a right corner $a=b=c=0$, and the assertion follows immediately from
\eqref{eq:content-formula} and \eqref{eq:polar-measure}.
\end{proof}

\begin{proposition}[Right and obtuse cases]\label{prop:nonacute-case}
If an orthocentric tetrahedron has a right or obtuse corner, then
\begin{equation}\label{eq:stronger-bound}
 \Sigma(T)\geq\frac{19\pi}{8}
 >6\arccos\frac13.
\end{equation}
\end{proposition}

\begin{proof}
Let the exceptional corner be number $4$.  Then $\beta_4\geq\pi/2$, while
the other three corners are acute.  Lemmas~\ref{lem:convexity} and
\ref{lem:obtuse-corner}, together with Theorem~\ref{thm:polar-sum}, give
\begin{align*}
 \sum_{i=1}^4\Omega_i
 &\geq \beta_4+\frac12(\beta_1+\beta_2+\beta_3)\\
 &=\beta_4+\frac12(\pi-\beta_4)
 =\frac{\pi+\beta_4}{2}
 \geq\frac{3\pi}{4}.
\end{align*}
Equation \eqref{eq:content-dihedral} yields
$\Sigma(T)\geq19\pi/8$.

For completeness, $19\pi/8>6\arccos(1/3)$ is equivalent to
$\cos(19\pi/48)<1/3$.  The latter follows from
\[
 \cos^2\frac{19\pi}{48}
 =\frac{1-\cos(5\pi/24)}{2}<\frac19.
\]
Indeed,
\[
 \cos^2\frac{5\pi}{24}
 =\frac12\left(1+\frac{\sqrt6-\sqrt2}{4}\right).
\]
Since $\sqrt3<7/4$, we have
$(\sqrt6-\sqrt2)^2=8-4\sqrt3>1$.  Hence
\[
 \cos^2\frac{5\pi}{24}>\frac58>\frac{49}{81},
\]
which proves $\cos(5\pi/24)>7/9$.
\end{proof}

Combining the two cases proves the central theorem.

\begin{theorem}[Sharp lower bound]\label{thm:lower-bound}
Every nondegenerate orthocentric tetrahedron satisfies
\begin{equation}\label{eq:main-lower-bound}
 \boxed{\ \Sigma(T)\geq6\arccos\frac13\ }.
\end{equation}
Equality holds if and only if $T$ is regular.
\end{theorem}

\section{Kites and the exact range}\label{sec:kites}

An equilateral-base kite has an equilateral base and three equal lateral
edges.  Normalize the base edge to $1$ and denote the lateral edge by $t$.
The nondegeneracy condition is
\begin{equation}\label{eq:kite-domain}
 t>\frac1{\sqrt3}.
\end{equation}
Every such kite is orthocentric.  In fact, in
\eqref{eq:orthocentric-metric} take the three base parameters equal to
$1/2$ and the apex parameter equal to $t^2-1/2$.

The three base-edge dihedral angles and the three lateral-edge dihedral
angles give the sum
\begin{equation}\label{eq:kite-sum}
 K(t)=3\left(
 \arccos\frac{1}{\sqrt{3(4t^2-1)}}
 +\arccos\frac{2t^2-1}{4t^2-1}
 \right).
\end{equation}
This formula follows directly by placing the apex above the center of the
equilateral base and taking scalar products of face normals.

Differentiation simplifies completely:
\begin{equation}\label{eq:kite-derivative}
 K'(t)=
 \frac{6(t-1)}{(4t^2-1)\sqrt{3t^2-1}}.
\end{equation}
Thus $K$ decreases on $(1/\sqrt3,1]$, reaches its unique minimum at the
regular tetrahedron $t=1$, and increases for $t>1$.  Moreover,
\begin{equation}\label{eq:kite-limits}
 K(1)=6\arccos\frac13,
 \qquad
 \lim_{t\downarrow1/\sqrt3}K(t)=3\pi.
\end{equation}
Consequently the short-kite branch alone realizes every value in
$[6\arccos(1/3),3\pi)$.

\begin{theorem}[Exact range]\label{thm:range}
The set of dihedral-angle sums of nondegenerate orthocentric tetrahedra is
\begin{equation}\label{eq:exact-range}
 \boxed{
 \left\{\Sigma(T):T\text{ orthocentric}\right\}
 =\left[6\arccos\frac13,\,3\pi\right).
 }
\end{equation}
\end{theorem}

\begin{proof}
Theorem~\ref{thm:lower-bound} gives the left endpoint and its equality case.
The universal upper bound in \eqref{eq:gaddum} gives $\Sigma(T)<3\pi$.
Finally, \eqref{eq:kite-sum}--\eqref{eq:kite-limits} show that the
orthocentric kite subfamily realizes every value between these bounds.
\end{proof}

The distinction between signed and positive parameters can now be
quantified exactly.  In the
kite parametrization above, the signed apex parameter is $a_4=t^2-1/2$.
The flattening branch $t\downarrow1/\sqrt3$ therefore enters the component
$a_4<0$ once $t<1/\sqrt2$.  It belongs to the full orthocentric family but
not to the positive Power-$2$ branch $\mathcal P_2^+$.

\begin{corollary}[Exact range for the positive Power-$2$ branch]
\label{cor:k2-range}
For tetrahedra in $\mathcal P_2^+$,
\begin{equation}\label{eq:k2-range}
 \boxed{
 \left\{\Sigma(T):T\in\mathcal P_2^+\right\}
 =\left[6\arccos\frac13,\,\frac{5\pi}{2}\right).
 }
\end{equation}
\end{corollary}

\begin{proof}
Every tetrahedron in $\mathcal P_2^+$ is orthocentric, so
Theorem~\ref{thm:lower-bound} supplies the lower bound and its equality case.
Its dihedral angles are nonobtuse; equivalently, this also follows from the
positive formula \eqref{eq:y-dihedral} below.  The sharp general bound for a
nonobtuse tetrahedron is
\[
 \Sigma(T)<\frac{5\pi}{2};
\]
see \cite[Theorem~3.1, $n=3$]{BrandtsCihangirKrizek2015}.

It remains only to realize every intermediate value.  Take the kite branch
$t\geq1$.  It lies in $\mathcal P_2^+$: take the three base parameters
$y_1=y_2=y_3=1/2$ and the apex parameter
\[
 y_4=t^2-\frac12>0.
\]
By \eqref{eq:kite-derivative}, $K(t)$ increases strictly from
$K(1)=6\arccos(1/3)$, and formula \eqref{eq:kite-sum} gives
\[
 \lim_{t\to\infty}K(t)
 =3\left(\frac\pi2+\frac\pi3\right)=\frac{5\pi}{2}.
\]
Thus this one-parameter positive Power-$2$ subfamily realizes the entire interval in
\eqref{eq:k2-range}.
\end{proof}

\section{\texorpdfstring{The Power-$k$ family}{The Power-k family}}\label{sec:power-family}

We now vary the exponent in the Power-$k$ family $\mathcal P_k$ of
Definition~\ref{def:power-k}.  The signed parameters $y_i$ are kept directly;
for even or nonintegral $k$, a negative $y_i$ need not be represented as a
real $k$th power.  We first record two simple geometric facts used at both
ends of the range.

\begin{lemma}[Kites occur for every exponent]\label{lem:all-k-kites}
For every $k\ne0$, $\mathcal P_k$ contains every equilateral-base
kite.
\end{lemma}

\begin{proof}
Take $y_1=y_2=y_3=u>0$ and $y_4=v$.  The base edge and the three lateral
edges are
\[
 b=(2u)^{1/k},\qquad \ell=(u+v)^{1/k}.
\]
Given an arbitrary ratio $t=\ell/b>0$, put
\begin{equation}\label{eq:kite-power-parameter}
 \frac vu=2t^k-1.
\end{equation}
Then $v>-u$, so all pair sums are positive and at most $v$ is negative.
For $t>1/\sqrt3$ the resulting kite is nondegenerate.  Hence all kites in
Section~\ref{sec:kites} occur.
\end{proof}

\begin{lemma}[Convex planar degeneration]\label{lem:convex-planar}
Suppose a continuous family of nondegenerate tetrahedra tends to four
distinct coplanar points in convex position.  Then the dihedral angles at
the two diagonals of the limiting quadrilateral tend to $\pi$, the other
four dihedral angles tend to $0$, and consequently
\begin{equation}\label{eq:convex-planar-limit}
 \Sigma(T)\longrightarrow2\pi.
\end{equation}
\end{lemma}

\begin{proof}
Orient the limiting plane horizontally.  Along a boundary edge, the two
incident faces approach the plane from opposite normal orientations, so the
internal dihedral angle tends to $0$.  Along a diagonal, the third vertices
lie on opposite sides of the diagonal in the limiting plane; the two face
normals therefore approach the same orientation and the internal dihedral
angle tends to $\pi$.  There are four boundary edges and two diagonals.
Equivalently, the statement follows by taking scalar products of the face
normals before passing to the limit.
\end{proof}

\subsection{\texorpdfstring{The exact range for $k<1$, $k\ne0$}
{The exact range for k below 1, excluding 0}}

\begin{theorem}\label{thm:k-below-one}
For every $k<1$, $k\ne0$, the dihedral-angle sums of Power-$k$
tetrahedra have the exact range
\begin{equation}\label{eq:k-below-one-range}
 \boxed{\{\Sigma(T):T\in\mathcal P_k\}=(2\pi,3\pi).}
\end{equation}
The lower endpoint is already approached with all four parameters positive.
\end{theorem}

\begin{proof}
The universal bounds \eqref{eq:gaddum} give containment in $(2\pi,3\pi)$.
Lemma~\ref{lem:all-k-kites} and the short-kite branch give every value in
$[6\arccos(1/3),3\pi)$.

It remains to reach the lower endpoint.  Put
\begin{equation}\label{eq:two-pair-path}
 (y_1,y_2,y_3,y_4)=(r,r,1,1),\qquad r>0,
\end{equation}
and set $q=2/k$.  The two distinguished opposite edges and the four equal
cross-edges are
\[
 a=(2r)^{1/k},\qquad c=2^{1/k},\qquad
 \ell=(1+r)^{1/k}.
\]
They are realized by
\begin{equation}\label{eq:two-pair-coordinates}
 \begin{aligned}
 A&=(-a/2,0,0),&B&=(a/2,0,0),\\
 C&=(0,h,c/2),&D&=(0,h,-c/2),
 \end{aligned}
 \qquad
 h=\frac12\sqrt{F_k(r)},
\end{equation}
where
\begin{equation}\label{eq:two-pair-feasibility}
 F_k(r)=4(1+r)^q-(2r)^q-2^q.
\end{equation}
At $r=1$ all six edges are equal and $F_k(1)>0$.  If $0<k<1$, then
$q>2$ and
\[
 F_k(0)=4-2^q<0.
\]
If $k<0$, then $q<0$ and $F_k(r)\to-\infty$ as $r\downarrow0$.
In either case, the component of $\{r:F_k(r)>0\}$ containing $r=1$ has a
lower endpoint $r_0\in(0,1)$ with $F_k(r_0)=0$.

As $r\downarrow r_0$, the height in
\eqref{eq:two-pair-coordinates} tends to zero.  The segments $AB$ and $CD$
remain nonzero, meet orthogonally at their common midpoint, and are the two
diagonals of the limiting convex quadrilateral.  Lemma~\ref{lem:convex-planar}
therefore gives
\[
 \Sigma(r)\longrightarrow2\pi,
 \qquad
 \Sigma(1)=6\arccos\frac13.
\]
The sum is continuous on $(r_0,1]$.  The intermediate value theorem supplies
every value in $(2\pi,6\arccos(1/3)]$; monotonicity is unnecessary.  Combining
this path with the kite path proves \eqref{eq:k-below-one-range}.
\end{proof}

\subsection{\texorpdfstring{The exact range for $k>2$}
{The exact range for k above 2}}

Fix $k>2$ and put $p=1/k\in(0,1/2)$.  Consider the signed Power-$k$ path
\begin{equation}\label{eq:signed-thin-path}
 (y_1,y_2,y_3,y_4)=(1,t,t,-rt),
 \qquad t>0,\qquad0\leq r<1.
\end{equation}
Set
\begin{equation}\label{eq:signed-thin-lengths}
 L=(1+t)^p,\quad M=(1-rt)^p,\quad
 a=(2t)^p,\quad b=((1-r)t)^p.
\end{equation}
The face on vertices $2,3,4$ is isosceles with sides $a,b,b$ and is
nondegenerate precisely when
\begin{equation}\label{eq:rc}
 r<r_c:=1-2^{1-k}.
\end{equation}
For $r<r_c$ use the coordinates
\begin{equation}\label{eq:signed-thin-coordinates}
 \begin{aligned}
 B&=(-a/2,0,0),& C&=(a/2,0,0),\\
 D&=(0,\eta,0),& A&=(0,z,h),
 \end{aligned}
\end{equation}
where
\begin{equation}\label{eq:signed-thin-height}
 \eta=\sqrt{b^2-a^2/4},\qquad
 z=\frac{L^2-M^2+b^2-a^2/2}{2\eta},\qquad
 h^2=L^2-a^2/4-z^2.
\end{equation}
These formulas impose $AB=AC=L$ and $AD=M$.

\begin{lemma}[Boundary layer of the signed path]\label{lem:signed-boundary}
Let $k>2$ and use \eqref{eq:signed-thin-path}--\eqref{eq:signed-thin-height}.
For every sufficiently small $t>0$ there is a first value
$r_*(t)\in(0,r_c)$ such that
\[
 h^2>0\quad(0\le r<r_*(t)),\qquad h(r_*(t))=0.
\]
Moreover, if $c_0=1-r_c=2^{1-k}$, then
\begin{equation}\label{eq:rstar-asymptotic}
 r_*(t)=r_c-\bigl(\lambda_*+o(1)\bigr)t^{2p},
 \qquad
 \lambda_*=
 \frac{2^{4p-6}}{2p\,c_0^{\,2p-1}}>0.
\end{equation}
At the limiting planar configuration one has $z<0<\eta$, so the four points
are in convex position.  Consequently
\begin{equation}\label{eq:signed-lower-limit}
 \Sigma(t,r)\longrightarrow2\pi
 \quad\text{as }r\uparrow r_*(t).
\end{equation}
\end{lemma}

\begin{proof}
At $r=0$,
\[
 \eta=t^p\sqrt{1-2^{2p-2}},\qquad z=O(t^p),\qquad h^2\longrightarrow1
 \quad(t\downarrow0),
\]
so the path begins in the nondegenerate region for small $t$.

Write $\delta=r_c-r$ and $c_0=1-r_c$.  Since
$c_0^{2p}=2^{2p-2}$, the first expression in
\eqref{eq:signed-thin-height} gives the exact identity
\begin{equation}\label{eq:eta-delta}
 \eta^2=t^{2p}\bigl((c_0+\delta)^{2p}-c_0^{2p}\bigr).
\end{equation}
Let
\[
 N(t,r)=L^2-M^2+b^2-a^2/2
\]
be the numerator of $z$.  In the critical layer
$\delta=\lambda t^{2p}$, with fixed $\lambda>0$, Taylor expansion yields
\begin{align}
 \eta
 &\sim t^{2p}\sqrt{2p\,c_0^{\,2p-1}\lambda},\label{eq:eta-layer}\\
 N(t,r)
 &\sim-2^{2p-2}t^{2p},\label{eq:n-layer}
\end{align}
because $L^2-M^2=O(t)=o(t^{2p})$ when $2p<1$.  Therefore
\begin{equation}\label{eq:z-layer}
 z\longrightarrow
 -\frac{2^{2p-3}}
 {\sqrt{2p\,c_0^{\,2p-1}\lambda}},
 \qquad
 h^2\longrightarrow1-\frac{\lambda_*}{\lambda}.
\end{equation}
In particular, the limiting sign of $h^2$ is positive for
$\lambda>\lambda_*$ and negative for $0<\lambda<\lambda_*$.

It remains to rule out an earlier loss of feasibility.  Suppose
$t_n\downarrow0$, $r_n<r_c$, and put
$\delta_n=r_c-r_n$.  If
$\delta_n/t_n^{2p}\to\infty$, then
\eqref{eq:eta-delta} and the expression for $N$ give
$z(t_n,r_n)\to0$, hence $h^2(t_n,r_n)\to1$.  If instead
$\delta_n/t_n^{2p}$ stays bounded away from zero and infinity, a subsequence
lies in a critical layer and \eqref{eq:z-layer} applies.  Consequently, for
any fixed $\lambda_+>\lambda_*$ and all sufficiently small $t$,
$h^2(t,r)>0$ throughout
$0\le r\le r_c-\lambda_+t^{2p}$, whereas for any
$0<\lambda_-<\lambda_*$,
\[
 h^2\bigl(t,r_c-\lambda_-t^{2p}\bigr)<0.
\]
By continuity there is a first zero $r_*(t)$ between these two values.
Letting $\lambda_\pm\to\lambda_*$ proves
\eqref{eq:rstar-asymptotic}.

At this first zero, \eqref{eq:z-layer} shows $z<0$ for small $t$, while
$\eta>0$ because $r_*(t)<r_c$.  In the plane $h=0$, the segment $AD$ therefore
crosses $BC$ at an interior point, so $A,B,D,C$ are four distinct points in
convex position and $AD,BC$ are their diagonals.  Lemma~\ref{lem:convex-planar}
then proves \eqref{eq:signed-lower-limit}.
\end{proof}

\begin{theorem}\label{thm:k-above-two}
For every $k>2$, the dihedral-angle sums of Power-$k$ tetrahedra have the
exact range
\begin{equation}\label{eq:k-above-two-range}
 \boxed{\{\Sigma(T):T\in\mathcal P_k\}=(2\pi,3\pi).}
\end{equation}
\end{theorem}

\begin{proof}
The universal bounds \eqref{eq:gaddum} give containment in $(2\pi,3\pi)$.
At $r=0$ the base triangle $BCD$ shrinks as $t\downarrow0$ while its shape
remains fixed, whereas the three edges from $A$ tend to $1$.  The three
base-edge dihedral angles tend to $\pi/2$, and the three lateral-edge
dihedral angles tend to the three angles of the base triangle.  Hence
\begin{equation}\label{eq:tall-limit}
 \Sigma(t,0)\longrightarrow\frac{5\pi}{2}
 \quad(t\downarrow0).
\end{equation}
Choose $t$ so small that
$\Sigma(t,0)>6\arccos(1/3)$.  By Lemma~\ref{lem:signed-boundary}, the same
connected feasible path has $\Sigma(t,r)\to2\pi$ as
$r\uparrow r_*(t)$.  Continuity therefore supplies every value from $2\pi$
through the regular value.  Lemma~\ref{lem:all-k-kites} supplies every value
from the regular value to $3\pi$.  The universal strict bounds complete the
proof.
\end{proof}

\subsection{\texorpdfstring{The remaining interval $1<k<2$}
{The remaining interval between 1 and 2}}

\begin{remark}[Status of the Power-$1$ case]\label{rem:power-one-status}
By the discussion after Definition~\ref{def:power-k}, $\mathcal P_1$ is
exactly the 4-ball (balloon) family.  Theorem~9 of
\cite{KorotovKrizek2026} states that its exact range is
$[6\arccos(1/3),3\pi)$.  Its kite subfamily rigorously realizes that whole
interval, so the only unresolved logical step is exclusion of values below
the regular value.  The cited proof treats an isosceles spherical triangle
and says that the nonsymmetric case follows by a ``tedious calculation,''
but the calculation or a certificate for the global maximization is not
given.  We do not use that lower bound in any proof below.
\end{remark}

Theorems~\ref{thm:k-below-one}, \ref{thm:range}, and
\ref{thm:k-above-two} suggest the following conjecture.

\begin{conjecture}[Power-$k$ window]\label{conj:power-window}
For every $1<k<2$, every Power-$k$ tetrahedron satisfies
\begin{equation}\label{eq:power-window-bound}
 \Sigma(T)\geq6\arccos\frac13,
\end{equation}
with equality only for the regular tetrahedron.  Consequently its exact
range is
\[
 \left[6\arccos\frac13,3\pi\right).
\]
\end{conjecture}

It is convenient to put $q=2/k$, so $1<q<2$, and to write the squared
edges as
\begin{equation}\label{eq:q-metric}
 D_{ij}=d_{ij}^2=(y_i+y_j)^q.
\end{equation}
The following observation removes a potential feasibility obstruction from
any attempt to deform the exponent down to the orthocentric value $q=1$.

\begin{proposition}[Downward feasibility]\label{prop:snowflake-feasibility}
Fix real parameters $y_1,\ldots,y_4$ with $y_i+y_j>0$ for $i\ne j$.  If,
for some $q_0>0$, the six numbers
\[
 d_{ij}(q_0)=(y_i+y_j)^{q_0/2}
\]
are the edge lengths of a nondegenerate Euclidean tetrahedron, then the same
is true for $d_{ij}(q)=(y_i+y_j)^{q/2}$ for every $0<q\le q_0$.
\end{proposition}

\begin{proof}
For $0<q<q_0$ put $\alpha=q/q_0\in(0,1)$.  Then
\[
 d_{ij}(q)=d_{ij}(q_0)^\alpha.
\]
Schoenberg's snowflake theorem says that a Euclidean metric raised to a
power $\alpha\in(0,1)$ is again Euclidean \cite{Schoenberg1938}.  Its strict
finite form says that for distinct points and nonzero coefficients $c_i$
with $\sum_i c_i=0$,
\[
 \sum_{i,j}c_ic_jd_{ij}(q_0)^{2\alpha}<0.
\]
Equivalently, the centered Gram matrix of the four new points is positive
definite on the three-dimensional zero-sum subspace.  Hence the four points
are affinely independent and form a nondegenerate tetrahedron.  The case
$q=q_0$ is the hypothesis.
\end{proof}

Whenever the metric is feasible, let $\Sigma_y(q)$ denote its dihedral-angle
sum.  Proposition~\ref{prop:snowflake-feasibility} makes the following
inequality meaningful on the whole interval from a given $q$ down to $1$.

\begin{conjecture}[Exponent monotonicity]\label{conj:q-monotonicity}
For every fixed admissible parameter vector $y$, the function
$q\mapsto\Sigma_y(q)$ is nondecreasing on every feasible subinterval of
$[1,2]$.
\end{conjecture}

\begin{proposition}[Conditional reduction]\label{prop:conditional-window}
Conjecture~\ref{conj:q-monotonicity} implies the Power-$k$ window
Conjecture~\ref{conj:power-window}.  It also supplies the missing lower-bound
step in Remark~\ref{rem:power-one-status}.
\end{proposition}

\begin{proof}
Let $1\le k<2$, put $q=2/k\in(1,2]$, and suppose that the corresponding
Power-$k$ metric is feasible.  Proposition~\ref{prop:snowflake-feasibility}
shows that the same parameters at $q=1$ define a nondegenerate
orthocentric tetrahedron.  Therefore Theorem~\ref{thm:lower-bound} and the
conjectured monotonicity give
\[
 \Sigma_y(q)\ge\Sigma_y(1)\ge6\arccos\frac13.
\]
If equality holds, the orthocentric equality case forces all four parameters
to be equal, and hence the original tetrahedron is regular.  Finally,
Lemma~\ref{lem:all-k-kites} and the short-kite branch realize every value
from the regular value up to $3\pi$, while the universal upper bound is
strict.
\end{proof}

The next proposition proves unconditionally that the proposed equality case
is a strict local minimum throughout the entire interval.

\begin{proposition}[Strict local minimality]\label{prop:k-local-min}
Let $q>0$.  Subject to a normalization of scale, the regular point
$y_1=y_2=y_3=y_4$ is a strict local minimum of the dihedral sum.  More
precisely, if $\sum_i v_i=0$, then
\begin{equation}\label{eq:regular-expansion}
 \Sigma_q(1+\varepsilon v_1,\ldots,1+\varepsilon v_4)
 =6\arccos\frac13
 +\frac{q^2}{12\sqrt2}\varepsilon^2\sum_{i=1}^4v_i^2
 +O(\varepsilon^3).
\end{equation}
\end{proposition}

\begin{proof}
Base the vertex Gram matrix at vertex $1$:
\[
 G_{ab}=\frac{D_{1,a+1}+D_{1,b+1}-D_{a+1,b+1}}2,
 \qquad1\leq a,b\leq3,
\]
where $D_{ii}=0$, and put $H=G^{-1}$ and ${\bf e}=(1,1,1)^T$.  The six
dihedral cosines are
\begin{equation}\label{eq:inverse-gram-dihedrals}
 \frac{(H{\bf e})_a}{\sqrt{H_{aa}{\bf e}^TH{\bf e}}}
 \quad(1\leq a\leq3),
 \qquad
 -\frac{H_{ab}}{\sqrt{H_{aa}H_{bb}}}
 \quad(1\leq a<b\leq3).
\end{equation}
Substitute \eqref{eq:q-metric} and differentiate at $y_i=1$.  Scale
invariance makes the all-ones direction null.  Permutation invariance makes
the Hessian a scalar on the irreducible subspace $\sum_i v_i=0$.  It is
therefore enough to compute along $v=(1,-1,0,0)$.

After dividing all squared edge lengths by the irrelevant common factor
$2^q$, one has
\[
 D_{12}=D_{34}=1,\qquad
 D_{13}=D_{14}=\left(1+\frac\varepsilon2\right)^q,\qquad
 D_{23}=D_{24}=\left(1-\frac\varepsilon2\right)^q.
\]
Expanding the inverse-Gram formulas \eqref{eq:inverse-gram-dihedrals}, the six
dihedral cosines, in a suitable order, are
\begin{align*}
 c_1&=\frac13+\left(\frac{5q^2}{27}-\frac q9\right)\varepsilon^2
       +O(\varepsilon^3),\\
 c_2=c_3&=\frac13-\frac{2q}{9}\varepsilon
       +\frac{q-q^2}{18}\varepsilon^2+O(\varepsilon^3),\\
 c_4=c_5&=\frac13+\frac{2q}{9}\varepsilon
       +\frac{q-q^2}{18}\varepsilon^2+O(\varepsilon^3),\\
 c_6&=\frac13-\frac{q^2+q}{9}\varepsilon^2+O(\varepsilon^3).
\end{align*}
Since
\[
 \arccos\left(\frac13+u\right)
 =\arccos\frac13-\frac{3}{2\sqrt2}u
  -\frac{9}{32\sqrt2}u^2+O(u^3),
\]
summing the six expansions cancels the linear terms and gives
\[
 \Sigma_q(1+\varepsilon,1-\varepsilon,1,1)
 =6\arccos\frac13+\frac{q^2}{6\sqrt2}\varepsilon^2
 +O(\varepsilon^3).
\]
Thus the second derivative in this direction is $q^2/(3\sqrt2)$.  Since
$\sum_i v_i^2=2$, permutation symmetry on $\sum_i v_i=0$ yields precisely
\eqref{eq:regular-expansion}.
\end{proof}

The two exterior constructions change character inside the open interval.
For the positive two-pair path \eqref{eq:two-pair-path}, if $1<q<2$, then
for every $r>0$,
\begin{equation}\label{eq:middle-two-pair-positive}
 (2r)^q+2^q=2^q(1+r^q)<4(1+r)^q.
\end{equation}
Thus this disphenoid path never reaches a convex planar boundary.  For the
signed path \eqref{eq:signed-thin-path}, now $p=1/k>1/2$.  At the face
boundary $r=r_c$, the numerator of $z$ is
\[
 N(t,r_c)=(1+t)^{2p}-(1-r_ct)^{2p}-2^{2p-2}t^{2p}.
\]
Here the first difference is $2p(1+r_c)t+O(t^2)$, whereas
$t^{2p}=o(t)$; hence $N(t,r_c)>0$ for small $t$.  This reverses the sign that
created the convex crossing in Lemma~\ref{lem:signed-boundary}, so that
particular boundary-layer proof no longer applies.  These observations rule
out two proposed counterexample mechanisms, but do not classify all boundary
strata and are not a global proof.

The remaining task has therefore been reduced to a clean analytic question:
prove Conjecture~\ref{conj:q-monotonicity}, or replace it by a boundary and
critical-point classification after a scale normalization.  Opposite-edge
ordering, Regge transformations, and elementary averaging do not by
themselves control the sum of six arccosines.  Any future proof using those
ideas still needs a displayed sign decomposition, a sequence of stated
inequalities, or a formally checkable symbolic certificate.

\section{The isodynamic family: the full universal range}
\label{sec:isodynamic}

A tetrahedron is \emph{isodynamic} precisely when its six edge lengths admit
positive parameters $p_1,p_2,p_3,p_4$ such that
\begin{equation}\label{eq:isodynamic-metric}
 d_{ij}=p_i p_j \qquad (i\ne j).
\end{equation}
Equivalently, the products of the lengths in the three pairs of opposite
edges are equal.  This is the multiplicative metric characterization in
\cite[Theorem~4]{HajjaHammoudeh2014}.  The following result settles the
range completely.

\begin{theorem}[Exact isodynamic range]\label{thm:isodynamic-range}
The set of dihedral-angle sums of nondegenerate isodynamic tetrahedra is
\begin{equation}\label{eq:isodynamic-range}
 \boxed{
 \left\{\Sigma(T):T\text{ isodynamic}\right\}=(2\pi,3\pi).
 }
\end{equation}
Neither endpoint is attained.
\end{theorem}

\begin{proof}
The universal inequalities \eqref{eq:gaddum} show at once that the range is
contained in $(2\pi,3\pi)$.  We prove the reverse inclusion by joining the
regular tetrahedron to each endpoint through explicit isodynamic
subfamilies.

For the upper endpoint, take
\[
 (p_1,p_2,p_3,p_4)=(s,1,1,1).
\]
The three edges among vertices $2,3,4$ have length $1$, and the remaining
three have length $s$.  Thus this is exactly the kite of
Section~\ref{sec:kites}, with $s>1/\sqrt3$.  On the branch
$1/\sqrt3<s\leq1$, formula \eqref{eq:kite-sum} realizes every value in
\[
 \left[6\arccos\frac13,3\pi\right).
\]

For the lower endpoint, take
\begin{equation}\label{eq:isodynamic-lower-path}
 (p_1,p_2,p_3,p_4)=(t,t,1,1)
\end{equation}
and label the vertices $A,B,C,D$ in this order.  Its edge lengths are
\[
 AB=t^2,\qquad CD=1,\qquad
 AC=AD=BC=BD=t.
\]
Put
\[
 h=\frac12\sqrt{4t^2-t^4-1}.
\]
The coordinates
\begin{equation}\label{eq:isodynamic-coordinates}
 \begin{aligned}
 A&=(-t^2/2,0,0),& B&=(t^2/2,0,0),\\
 C&=(0,h,1/2),& D&=(0,h,-1/2)
 \end{aligned}
\end{equation}
realize the required lengths whenever
\begin{equation}\label{eq:isodynamic-domain}
 \sqrt{2-\sqrt3}<t<\sqrt{2+\sqrt3}.
\end{equation}
There are three types of internal dihedral angles: one at $AB$, one at
$CD$, and four equal angles at the cross-edges.  Direct scalar products of
the inward face normals in \eqref{eq:isodynamic-coordinates} give
\begin{align}
 \cos\theta_{AB}
 &=\frac{4t^2-t^4-2}{4t^2-t^4},\label{eq:iso-cos-ab}\\
 \cos\theta_{CD}
 &=\frac{4t^2-2t^4-1}{4t^2-1},\label{eq:iso-cos-cd}\\
 \cos\theta_{\times}
 &=\frac{t}{\sqrt{(4-t^2)(4t^2-1)}}.\label{eq:iso-cos-cross}
\end{align}
Hence
\begin{equation}\label{eq:isodynamic-path-sum}
 L(t)=\theta_{AB}(t)+\theta_{CD}(t)+4\theta_{\times}(t)
\end{equation}
is continuous throughout \eqref{eq:isodynamic-domain}.  At $t=1$ the
tetrahedron is regular, so $L(1)=6\arccos(1/3)$.  If
$t_0=\sqrt{2+\sqrt3}$, then \eqref{eq:iso-cos-ab}--\eqref{eq:iso-cos-cross}
give
\[
 \theta_{AB}\longrightarrow\pi,\qquad
 \theta_{CD}\longrightarrow\pi,\qquad
 \theta_{\times}\longrightarrow0
 \quad\text{as }t\uparrow t_0.
\]
Therefore $L(t)\to2\pi$.  The intermediate value theorem shows that this
path realizes every value in
$(2\pi,6\arccos(1/3)]$.  Monotonicity is not needed: for any number strictly
between the limit and the value at $t=1$, continuity on a sufficiently long
closed subinterval already supplies a preimage.  Together with the kite
branch, this proves that every number in $(2\pi,3\pi)$ occurs.  The strict
universal inequalities show that the endpoints do not occur.
\end{proof}

\begin{remark}
Thus the isodynamic theorem really does reduce to ``examples at the two
boundaries,'' but only after three points are checked carefully: both paths
must remain inside the isodynamic family, their tetrahedra must remain
nondegenerate before the limiting parameter, and the dihedral sum must vary
continuously.  Equations \eqref{eq:isodynamic-metric},
\eqref{eq:isodynamic-coordinates}, and \eqref{eq:isodynamic-path-sum} supply
exactly these checks.
\end{remark}

\section{A comparison family: isogonic tetrahedra}
\label{sec:isogonic}

The isogonic family is included here for comparison, not as another member of
the Power-$k$ chain.  The exact formulas below are proved, but the regular
lower bound remains conditional on Conjecture~\ref{conj:isogonic-reduction}.
This section should therefore be read as a self-contained open problem rather
than as part of the completed Power-$k$ classification.

A tetrahedron is \emph{isogonic} precisely when its squared edge lengths can
be written
\begin{equation}\label{eq:isogonic-metric}
 d_{ij}^2=p_i^2+p_i p_j+p_j^2,
 \qquad p_1,p_2,p_3,p_4>0;
\end{equation}
see \cite[Theorem~4]{HajjaHammoudeh2014}.  In contrast with the
orthocentric case, the following calculation exposes the whole dihedral sum
without a choice of branches or auxiliary solid-angle measures.

For complementary indices $\{i,j,k,l\}=\{1,2,3,4\}$, set
\begin{equation}\label{eq:isogonic-notation}
 q_{ij}=p_i^2+p_i p_j+p_j^2,
 \quad
 E=\sum_{1\leq r<s\leq4}p_rp_s,
 \quad
 h_i=p_jp_k+p_jp_l+p_kp_l.
\end{equation}

\begin{proposition}[Isogonic half-angle formula]
\label{prop:isogonic-formula}
Let $\theta_{ij}$ be the internal dihedral angle between the faces opposite
vertices $i$ and $j$; equivalently, it is the angle at edge $kl$.  Then
\begin{align}
 \sin^2\frac{\theta_{ij}}2
 &=\frac{p_i p_j q_{kl}}{h_i h_j},
 \label{eq:isogonic-sin-half}\\
 \tan^2\frac{\theta_{ij}}2
 &=\frac{p_i p_j q_{kl}}{p_kp_l E}.
 \label{eq:isogonic-tan-half}
\end{align}
Consequently
\begin{equation}\label{eq:isogonic-sum}
 \boxed{
 \Sigma_{\rm igo}(p_1,p_2,p_3,p_4)
 =2\sum_{1\leq i<j\leq4}
 \arctan\sqrt{\frac{p_i p_j q_{kl}}{p_kp_lE}} .
 }
\end{equation}
\end{proposition}

\begin{proof}
Heron's formula applied to the face opposite vertex $i$ simplifies to
\begin{equation}\label{eq:isogonic-face-area}
 A_i=\frac{\sqrt3}{4}h_i.
\end{equation}
The determinant of the edge Gram matrix based at any vertex simplifies to
\begin{equation}\label{eq:isogonic-volume}
 36V^2=\frac94p_1p_2p_3p_4E,
 \qquad
 V^2=\frac1{16}p_1p_2p_3p_4E.
\end{equation}
A direct cofactor calculation for the angle between two face normals then
gives
\begin{equation}\label{eq:isogonic-cos}
 \cos\theta_{ij}
 =1-\frac{2p_ip_jq_{kl}}{h_ih_j}.
\end{equation}
Using $\sin^2(\theta/2)=(1-\cos\theta)/2$ proves
\eqref{eq:isogonic-sin-half}.  Finally, direct expansion gives the useful
identity
\begin{equation}\label{eq:isogonic-identity}
 h_ih_j-p_ip_jq_{kl}=p_kp_lE.
\end{equation}
Dividing \eqref{eq:isogonic-sin-half} by its complementary cosine-square
proves \eqref{eq:isogonic-tan-half}, and summing the six principal
half-angles proves \eqref{eq:isogonic-sum}.
\end{proof}

The formula immediately explains why kites are central here.  If
$(p_1,p_2,p_3,p_4)=(a,b,b,b)$, then the base edges have squared length
$3b^2$ and the three lateral edges have squared length $a^2+ab+b^2$.
Writing $s=a/b$, the squared lateral-to-base ratio is
\begin{equation}\label{eq:isogonic-kite-ratio}
 t^2=\frac{s^2+s+1}{3}.
\end{equation}
As $s$ runs over $(0,\infty)$, $t$ runs over $(1/\sqrt3,\infty)$.
Therefore the isogonic family contains every equilateral-base kite, and in
particular it realizes
\begin{equation}\label{eq:isogonic-known-half}
 \left[6\arccos\frac13,3\pi\right).
\end{equation}
The unresolved direction is to prove that no isogonic tetrahedron lies below
the regular value.

Formula \eqref{eq:isogonic-sum} suggests the following precise reduction.

\begin{conjecture}[Ordered kite-reduction inequality]
\label{conj:isogonic-reduction}
Suppose $0<p_1\leq p_2\leq p_3\leq p_4$ and put
\[
 m=\frac{p_1+p_2+p_3}{3}.
\]
Then
\begin{equation}\label{eq:isogonic-reduction}
 \Sigma_{\rm igo}(p_1,p_2,p_3,p_4)
 \geq \Sigma_{\rm igo}(m,m,m,p_4).
\end{equation}
Equality should hold only when $p_1=p_2=p_3$.
\end{conjecture}

\begin{theorem}[Conditional exact isogonic range]
\label{thm:conditional-isogonic}
If Conjecture~\ref{conj:isogonic-reduction} holds, then
\begin{equation}\label{eq:conditional-isogonic-range}
 \left\{\Sigma(T):T\text{ isogonic}\right\}
 =\left[6\arccos\frac13,3\pi\right),
\end{equation}
and equality at the lower endpoint occurs only for the regular tetrahedron.
\end{theorem}

\begin{proof}
After relabelling, the parameters may be ordered.  The right-hand side of
\eqref{eq:isogonic-reduction} is the sum of an equilateral-base kite and is
therefore at least its regular value by \eqref{eq:kite-derivative}.  Equality
in the conjecture first forces $p_1=p_2=p_3$; equality in the kite bound then
forces $p_4=m$, so all four parameters are equal and the tetrahedron is
regular.  The upper bound is universal, and \eqref{eq:isogonic-known-half}
already supplies every value from the regular value to $3\pi$.
\end{proof}

The ordering in Conjecture~\ref{conj:isogonic-reduction} is essential to its
formulation.  Neither the half-angle formula nor ordinary pairwise averaging
provides the required sign.  We therefore record the reduction only as an
open problem: a proof must exhibit a readable chain of inequalities or a
formally checkable symbolic certificate.

\section{Discussion: why the kite coincidence is useful}

The equality of the orthocentric range and the kite range has
two distinct consequences.

First, because kites are a subfamily, their one-parameter calculation proves
sharpness, the unattained upper endpoint, and the realization of all
intermediate values.  It does \emph{not} prove the lower bound for the larger
family; that logical direction is supplied by Theorem~\ref{thm:lower-bound}.

Second, the coincidence suggests a stronger comparison statement: perhaps
each orthocentric tetrahedron can be symmetrized to a kite without increasing
its dihedral sum.  The signed metric parametrization
$d_{ij}^2=a_i+a_j$ makes this plausible.  If all $a_i\ne0$ and
\[
 S=\sum_{i=1}^4\frac1{a_i},
\]
then an inverse-Gram calculation gives, for the dihedral angle between the
faces opposite vertices $i$ and $j$,
\begin{equation}\label{eq:direct-dihedral}
 \cos\theta_{ij}
 =\frac{\operatorname{sgn}(Sa_ia_j)}
 {\sqrt{(Sa_i-1)(Sa_j-1)}}.
\end{equation}
When all $a_i$ are positive, the variables
\[
 y_i=\frac1{Sa_i},
 \qquad y_i>0,
 \qquad \sum_{i=1}^4y_i=1,
\]
reduce this to the symmetric formula
\begin{equation}\label{eq:y-dihedral}
 \cos\theta_{ij}
 =\sqrt{\frac{y_i y_j}{(1-y_i)(1-y_j)}}.
\end{equation}
Right-corner cases, for which one additive parameter vanishes, are obtained
from these formulas by a limiting argument.  A direct symmetrization theorem
would be stronger than the range result, but it is unnecessary here: the
constant polar sum converts the lower-bound problem into Jensen's inequality
for the one-variable function $f$.

The proof also separates the potentially troublesome obtuse component rather
efficiently.  Instead of carrying the branch signs through a lengthy
six-dihedral calculation, Lemma~\ref{lem:obtuse-corner} gives
$\Omega\geq\beta$ at the exceptional corner, while convexity gives the simple
supporting-line estimate $\Omega\geq\beta/2$ at each acute corner.  This yields
the stronger non-acute bound $19\pi/8$ and leaves equality possible only in
the all-acute component.

The Power-$k$ family exhibits two rigorously different exterior regimes.
The regular lower bound is proved here at $k=2$, while explicit convex-planar
degenerations lower the infimum to $2\pi$ for $k<1$ and for $k>2$.  In the
first regime the degeneration already occurs in the positive branch; in the
second it requires one negative Power-$k$ parameter.  The interval
$1<k<2$ remains open, as does the unexpanded global step at $k=1$ noted in
Remark~\ref{rem:power-one-status}.  Proposition~\ref{prop:snowflake-feasibility}
and Conjecture~\ref{conj:q-monotonicity} isolate a unified route to both
questions.

The families studied here illustrate several genuinely different proof
mechanisms.
For orthocentric tetrahedra, a constant sum for a secondary angle measure
creates the decisive convexity argument.  For isodynamic tetrahedra, the
universal theorem plus explicit endpoint paths settles the range, even
though the sum need not be monotone along the lower path.  For isogonic
tetrahedra, exact formulas and the kite subfamily settle sharpness and the
upper part of the range, but they do not yet exclude a hidden value below the
regular one.  Conjecture~\ref{conj:isogonic-reduction} isolates that missing
step in a form suitable for either a human factorization or a formally
checkable symbolic certificate.

\section*{Acknowledgements}
The author thanks Mowaffaq Hajja for his support and encouragement.
 AI tools assisted with limited editorial and computational work; they are
 not coauthors.  The author bears full responsibility for the mathematical
 content and final wording of this article.

\begin{thebibliography}{99}
\small
\setlength{\itemsep}{4pt}

\bibitem{BrandtsCihangirKrizek2015}
J.~Brandts, A.~Cihangir, and M.~K\v{r}\'\i\v{z}ek,
\newblock Tight bounds on angle sums of nonobtuse simplices,
\newblock \emph{Appl. Math. Comput.} \textbf{267} (2015), 397--408,
\newblock \href{https://doi.org/10.1016/j.amc.2015.02.035}
{doi:10.1016/j.amc.2015.02.035}.

\bibitem{EdmondsHajjaMartini2005}
A.~L. Edmonds, M.~Hajja, and H.~Martini,
\newblock Orthocentric simplices and their centers,
\newblock \emph{Results Math.} \textbf{47} (2005), 266--295;
\newblock \href{https://doi.org/10.1007/BF03323029}
{doi:10.1007/BF03323029}.

\bibitem{Gaddum1952}
J.~W. Gaddum,
\newblock The sums of the dihedral and trihedral angles in a tetrahedron,
\newblock \emph{Amer. Math. Monthly} \textbf{59} (1952), 370--371,
\newblock \href{https://doi.org/10.1080/00029890.1952.11988143}
{doi:10.1080/00029890.1952.11988143}.

\bibitem{HajjaHammoudeh2014}
M.~Hajja and I.~Hammoudeh,
\newblock The sum of measures of the angles of a simplex,
\newblock \emph{Beitr. Algebra Geom.} \textbf{55} (2014), 453--470,
\newblock \href{https://doi.org/10.1007/s13366-013-0160-8}
{doi:10.1007/s13366-013-0160-8}.

\bibitem{Jarden1941}
D.~Jarden (D.~Juzuk),
\newblock On the sum of the dihedral and trihedral angles in a tetrahedron
(in Hebrew),
\newblock \emph{Technika Umada (Technics and Science)} \textbf{4} (1941),
no.~5 (26), 8--9.

\bibitem{Katsuura2021}
H.~Katsuura,
\newblock Dihedral angles of 4-ball tetrahedra,
\newblock \emph{J. Geom. Graph.} \textbf{25} (2021), no.~2, 197--204.

\bibitem{KorotovKrizek2026}
S.~Korotov and M.~K\v{r}\'\i\v{z}ek,
\newblock Two-sided bounds for dihedral angle sums of path and 4-ball
tetrahedra,
\newblock arXiv:2601.08304 [math.MG] (2026),
\newblock \href{https://doi.org/10.48550/arXiv.2601.08304}
{doi:10.48550/arXiv.2601.08304}.

\bibitem{Schoenberg1938}
I.~J. Schoenberg,
\newblock Metric spaces and positive definite functions,
\newblock \emph{Trans. Amer. Math. Soc.} \textbf{44} (1938), 522--536,
\newblock \href{https://doi.org/10.1090/S0002-9947-1938-1501980-0}
{doi:10.1090/S0002-9947-1938-1501980-0}.

\end{thebibliography}

\end{document}
